
\documentclass[12pt, a4paper, openany, twoside]{book}
\usepackage[left=1in,right=1in,bottom=0.5in,top=0.8in]{geometry}
\usepackage{amsfonts}
\usepackage{amsmath, amssymb}
\usepackage{graphicx}
\usepackage[font=small,labelfont=bf]{caption}
\usepackage{epstopdf}
\usepackage[pdfpagelabels,hyperindex]{hyperref}
\usepackage{xcolor}
\usepackage{amsthm}
\usepackage{float}
\usepackage{pgfplots}
\usepackage{listings}
\usepackage{longtable}
\usepackage{mathrsfs}
\usepackage{mathtools}
\usepackage{hyperref}
\usepackage[most]{tcolorbox}
\usepackage{amssymb}
\usepackage{algorithm}
\usepackage[noend]{algpseudocode}
\usepackage{tikz-cd}
\usepackage{comment}
\usepackage{mathpartir}

\linespread{1.3} 

\hypersetup{
pdftitle={.},
pdfauthor={Hojune Lee},
}

\newcommand{\overbar}[1]{\mkern 1.5mu\overline{\mkern-1.5mu#1\mkern-1.5mu}\mkern 1.5mu}
\DeclareRobustCommand{\stirling}{\genfrac\{\}{0pt}{}}
\newtheorem{thm}{Theorem}[section]
\newtheorem{cor}[thm]{Corollary}
\newtheorem{prop}[thm]{Proposition}
\newtheorem{lem}[thm]{Lemma}
\newtheorem{conj}[thm]{Conjecture}
\newtheorem{quest}[thm]{Question}
\newtheorem{ppty}[thm]{Property}
\newtheorem{ppties}[thm]{Properties}
\newtheorem{axiom}[thm]{Axiom}
\newtheorem{claim}[thm]{Claim}
\newtheorem{prob}[thm]{Problem}


\theoremstyle{definition}
\newtheorem{defn}[thm]{Definition}
\newtheorem{defns}[thm]{Definitions}
\newtheorem{con}[thm]{Construction}
\newtheorem{exmp}[thm]{Example}
\newtheorem{exmps}[thm]{Examples}
\newtheorem{notn}[thm]{Notation}
\newtheorem{notns}[thm]{Notations}
\newtheorem{addm}[thm]{Addendum}
\newtheorem{exer}[thm]{Exercise}
\newtheorem{limit}[thm]{Limitation}


\theoremstyle{remark}
\newtheorem{rem}[thm]{Remark}
\newtheorem{rems}[thm]{Remarks}
\newtheorem{warn}[thm]{Warning}
\newtheorem{sch}[thm]{Scholium}
\newenvironment{thmbox}[1][]{\begin{tcolorbox}[colback=yellow!10!white,colframe=red!75!black,title={#1}]}{\end{tcolorbox}}


% newcommand bb
    \newcommand{\BA}{{\mathbb {A}}} \newcommand{\BB}{{\mathbb {B}}}
    \newcommand{\BC}{{\mathbb {C}}} \newcommand{\BD}{{\mathbb {D}}}
    \newcommand{\BE}{{\mathbb {E}}} \newcommand{\BF}{{\mathbb {F}}}
    \newcommand{\BG}{{\mathbb {G}}} \newcommand{\BH}{{\mathbb {H}}}
    \newcommand{\BI}{{\mathbb {I}}} \newcommand{\BJ}{{\mathbb {J}}}
    \newcommand{\BK}{{\mathbb {U}}} \newcommand{\BL}{{\mathbb {L}}}
    \newcommand{\BM}{{\mathbb {M}}} \newcommand{\BN}{{\mathbb {N}}}
    \newcommand{\BO}{{\mathbb {O}}} \newcommand{\BP}{{\mathbb {P}}}
    \newcommand{\BQ}{{\mathbb {Q}}} \newcommand{\BR}{{\mathbb {R}}}
    \newcommand{\BS}{{\mathbb {S}}} \newcommand{\BT}{{\mathbb {T}}}
    \newcommand{\BU}{{\mathbb {U}}} \newcommand{\BV}{{\mathbb {V}}}
    \newcommand{\BW}{{\mathbb {W}}} \newcommand{\BX}{{\mathbb {X}}}
    \newcommand{\BY}{{\mathbb {Y}}} \newcommand{\BZ}{{\mathbb {Z}}}

% newcommand  scr
    \newcommand{\sA}{{\mathscr {A}}} \newcommand{\sB}{{\mathscr {B}}}
    \newcommand{\sC}{{\mathscr {C}}} \newcommand{\sD}{{\mathscr {D}}}
    \newcommand{\sE}{{\mathscr {E}}} \newcommand{\sF}{{\mathscr {F}}}
    \newcommand{\sG}{{\mathscr {G}}} \newcommand{\sH}{{\mathscr {H}}}
    \newcommand{\sI}{{\mathscr {I}}} \newcommand{\sJ}{{\mathscr {J}}}
    \newcommand{\sK}{{\mathscr {K}}} \newcommand{\sL}{{\mathscr {L}}}
    \newcommand{\sN}{{\mathscr {N}}} \newcommand{\sM}{{\mathscr {M}}}
    \newcommand{\sO}{{\mathscr {O}}} \newcommand{\sP}{{\mathscr {P}}}
    \newcommand{\sQ}{{\mathscr {Q}}} \newcommand{\sR}{{\mathscr {R}}}
    \newcommand{\sS}{{\mathscr {S}}} \newcommand{\sT}{{\mathscr {T}}}
    \newcommand{\sU}{{\mathscr {U}}} \newcommand{\sV}{{\mathscr {V}}}
    \newcommand{\sW}{{\mathscr {W}}} \newcommand{\sX}{{\mathscr {X}}}
    \newcommand{\sY}{{\mathscr {Y}}} \newcommand{\sZ}{{\mathscr {Z}}}


% newcommand cal
    \newcommand{\CA}{{\mathcal {A}}} \newcommand{\CB}{{\mathcal {B}}}
    \newcommand{\CC}{{\mathcal {C}}} \newcommand{\CD}{{\mathcal {D}}}
    \newcommand{\CE}{{\mathcal {E}}} \newcommand{\CF}{{\mathcal {F}}}
    \newcommand{\CG}{{\mathcal {G}}} \newcommand{\CH}{{\mathcal {H}}}
    \newcommand{\CI}{{\mathcal {I}}} \newcommand{\CJ}{{\mathcal {J}}}
    \newcommand{\CK}{{\mathcal {K}}} \newcommand{\CL}{{\mathcal {L}}}
    \newcommand{\CM}{{\mathcal {M}}} \newcommand{\CN}{{\mathcal {N}}}
    \newcommand{\CO}{{\mathcal {O}}} \newcommand{\CP}{{\mathcal {P}}}
    \newcommand{\CQ}{{\mathcal {Q}}} \newcommand{\CR}{{\mathcal {R}}}
    \newcommand{\CS}{{\mathcal {S}}} \newcommand{\CT}{{\mathcal {T}}}
    \newcommand{\CU}{{\mathcal {U}}} \newcommand{\CV}{{\mathcal {V}}}
    \newcommand{\CW}{{\mathcal {W}}} \newcommand{\CX}{{\mathcal {X}}}
    \newcommand{\CY}{{\mathcal {Y}}} \newcommand{\CZ}{{\mathcal {Z}}}

    % newcommand frak
     \newcommand{\fa}{{\mathfrak{a}}}  \newcommand{\fb}{{\mathfrak{b}}}
     \newcommand{\fc}{{\mathfrak{c}}}  \newcommand{\fd}{{\mathfrak{d}}}
     \newcommand{\fe}{{\mathfrak{e}}}  \newcommand{\ff}{{\mathfrak{f}}}
     \newcommand{\fg}{{\mathfrak{g}}}  \newcommand{\fh}{{\mathfrak{h}}}
     \newcommand{\fii}{{\mathfrak{i}}}  \newcommand{\fj}{{\mathfrak{j}}}
     \newcommand{\fk}{{\mathfrak{m}}}  \newcommand{\fl}{{\mathfrak{l}}}
     \newcommand{\fm}{{\mathfrak{m}}}  \newcommand{\fn}{{\mathfrak{n}}}
     \newcommand{\fo}{{\mathfrak{o}}}  \newcommand{\fp}{{\mathfrak{p}}}
     \newcommand{\fq}{{\mathfrak{q}}}  \newcommand{\fr}{{\mathfrak{r}}}
     \newcommand{\fs}{{\mathfrak{s}}}  \newcommand{\ft}{{\mathfrak{t}}}
     \newcommand{\fu}{{\mathfrak{u}}}  \newcommand{\fv}{{\mathfrak{v}}}
     \newcommand{\fw}{{\mathfrak{w}}}  \newcommand{\fx}{{\mathfrak{x}}}
     \newcommand{\fy}{{\mathfrak{y}}}  \newcommand{\fz}{{\mathfrak{z}}}

    \newcommand{\fA}{{\mathfrak{A}}}  \newcommand{\fB}{{\mathfrak{B}}}
     \newcommand{\fC}{{\mathfrak{C}}}  \newcommand{\fD}{{\mathfrak{D}}}
     \newcommand{\fE}{{\mathfrak{E}}}  \newcommand{\fF}{{\mathfrak{F}}}
     \newcommand{\fG}{{\mathfrak{G}}}  \newcommand{\fH}{{\mathfrak{H}}}
     \newcommand{\fI}{{\mathfrak{I}}}  \newcommand{\fJ}{{\mathfrak{J}}}
     \newcommand{\fK}{{\mathfrak{K}}}  \newcommand{\fL}{{\mathfrak{L}}}
     \newcommand{\fM}{{\mathfrak{M}}}  \newcommand{\fN}{{\mathfrak{N}}}
     \newcommand{\fO}{{\mathfrak{O}}}  \newcommand{\fP}{{\mathfrak{P}}}
     \newcommand{\fQ}{{\mathfrak{Q}}}  \newcommand{\fR}{{\mathfrak{R}}}
     \newcommand{\fS}{{\mathfrak{S}}}  \newcommand{\fT}{{\mathfrak{T}}}
     \newcommand{\fU}{{\mathfrak{U}}}  \newcommand{\fV}{{\mathfrak{V}}}
     \newcommand{\fW}{{\mathfrak{W}}}  \newcommand{\fX}{{\mathfrak{X}}}
     \newcommand{\fY}{{\mathfrak{Y}}}  \newcommand{\fZ}{{\mathfrak{Z}}}



 % newcommand :rm
     \newcommand{\RA}{{\mathrm {A}}} \newcommand{\RB}{{\mathrm {B}}}
    \newcommand{\RC}{{\mathrm {C}}} \newcommand{\RD}{{\mathrm {D}}}
    \newcommand{\RE}{{\mathrm {E}}} \newcommand{\RF}{{\mathrm {F}}}
    \newcommand{\RG}{{\mathrm {G}}} \newcommand{\RH}{{\mathrm {H}}}
    \newcommand{\RI}{{\mathrm {I}}} \newcommand{\RJ}{{\mathrm {J}}}
    \newcommand{\RK}{{\mathrm {K}}} \newcommand{\RL}{{\mathrm {L}}}
    \newcommand{\RM}{{\mathrm {M}}} \newcommand{\RN}{{\mathrm {N}}}
    \newcommand{\RO}{{\mathrm {O}}} \newcommand{\RP}{{\mathrm {P}}}
    \newcommand{\RQ}{{\mathrm {Q}}} \newcommand{\RR}{{\mathrm {R}}}
    \newcommand{\RS}{{\mathrm {S}}} \newcommand{\RT}{{\mathrm {T}}}
    \newcommand{\RU}{{\mathrm {U}}} \newcommand{\RV}{{\mathrm {V}}}
    \newcommand{\RW}{{\mathrm {W}}} \newcommand{\RX}{{\mathrm {X}}}
    \newcommand{\RY}{{\mathrm {Y}}} \newcommand{\RZ}{{\mathrm {Z}}}

    \newcommand{\Ad}{{\mathrm{Ad}}} \newcommand{\Aut}{{\mathrm{Aut}}}
    \newcommand{\Br}{{\mathrm{Br}}} \newcommand{\Ch}{{\mathrm{Ch}}}
    \newcommand{\cod}{{\mathrm{cod}}} \newcommand{\cont}{{\mathrm{cont}}}
    \newcommand{\cl}{{\mathrm{cl}}}   \newcommand{\Cl}{{\mathrm{Cl}}}
    \newcommand{\disc}{{\mathrm{disc}}}\newcommand{\Eis}{{\mathrm{Eis}}}
    \newcommand{\Div}{{\mathrm{Div}}} \renewcommand{\div}{{\mathrm{div}}}
    \newcommand{\End}{{\mathrm{End}}} \newcommand{\Frob}{{\mathrm{Frob}}}
    \newcommand{\Gal}{{\mathrm{Gal}}} \newcommand{\GL}{{\mathrm{GL}}}
    \newcommand{\Hom}{{\mathrm{Hom}}} \renewcommand{\Im}{{\mathrm{Im}}}
    \newcommand{\Ind}{{\mathrm{Ind}}} \newcommand{\ind}{{\mathrm{ind}}}
    \newcommand{\inv}{{\mathrm{inv}}}
    \newcommand{\Isom}{{\mathrm{Isom}}} \newcommand{\Jac}{{\mathrm{Jac}}}
    \newcommand{\ad}{{\mathrm{ad}}}  \newcommand{\Tr}{{\mathrm{Tr}}}
    \newcommand{\Ker}{{\mathrm{Ker}}} \newcommand{\Ros}{{\mathrm{Ros}}}
    \newcommand{\Lie}{{\mathrm{Lie}}} \newcommand{\Hol}{{\mathrm{Hol}}}

    \newcommand{\cyc}{{\mathrm{cyc}}}\newcommand{\id}{{\mathrm{id}}}
    \newcommand{\new}{{\mathrm{new}}} \newcommand{\NS}{{\mathrm{NS}}}
    \newcommand{\ord}{{\mathrm{ord}}} \newcommand{\rank}{{\mathrm{rank}}}
    \newcommand{\PGL}{{\mathrm{PGL}}} \newcommand{\Pic}{\mathrm{Pic}}
    \newcommand{\cond}{\mathrm{cond}} \newcommand{\Is}{{\mathrm{Is}}}
    \renewcommand{\Re}{{\mathrm{Re}}} \newcommand{\reg}{{\mathrm{reg}}}
    \newcommand{\Res}{{\mathrm{Res}}} \newcommand{\Sel}{{\mathrm{Sel}}}
    \newcommand{\RTr}{{\mathrm{Tr}}} \newcommand{\alg}{{\mathrm{alg}}}
    \newcommand{\PSL}{{\mathrm{PSL}}}

\newcommand{\coker}{{\mathrm{coker}}}
\newcommand{\val}{{\mathrm{val}}} \newcommand{\sign}{{\mathrm{sign}}}
\newcommand{\mult}{{\mathrm{mult}}} \newcommand{\Vol}{{\mathrm{Vol}}}
\newcommand{\Meas}{{\mathrm{Meas}}}\renewcommand{\mod}{\ \mathrm{mod}\ }
\newcommand{\Ann}{\mathrm{Ann}}
\newcommand{\Tor}{\mathrm{Tor}}
\newcommand{\Supp}{\mathrm{Supp}}\newcommand{\supp}{\mathrm{supp}}
\newcommand{\Max}{\mathrm{Max}}
\newcommand{\Coker}{\mathrm{Coker}}
\newcommand{\Stab}{\mathrm{Stab}}
\newcommand{\Irr}{\mathrm{Irr}}\newcommand{\Inf}{\mathrm{Inf}}\newcommand{\Sup}{\mathrm{Sup}}
\newcommand{\rk}{\mathrm{rk}}\newcommand{\Fil}{\mathrm{Fil}}
\newcommand{\Sim}{{\mathrm{Sim}}} \newcommand{\SL}{{\mathrm{SL}}}
\newcommand{\Spec}{{\mathrm{Spec}}} \newcommand{\SO}{{\mathrm{SO}}}
\newcommand{\SU}{{\mathrm{SU}}} \newcommand{\Sym}{{\mathrm{Sym}}}
\newcommand{\sgn}{{\mathrm{sgn}}} \newcommand{\tr}{{\mathrm{tr}}}
\newcommand{\tor}{{\mathrm{tor}}}  \newcommand{\ur}{{\mathrm{ur}}}
\newcommand{\vol}{{\mathrm{vol}}}  \newcommand{\ab}{{\mathrm{ab}}}
\newcommand{\Sh}{{\mathrm{Sh}}} \newcommand{\Ell}{{\mathrm{Ell}}}
\newcommand{\Char}{{\mathrm{Char}}}\newcommand{\Tate}{{\mathrm{Tate}}}
\newcommand{\corank}{{\mathrm{corank}}} \newcommand{\Cond}{{\mathrm{Cond}}}
\newcommand{\Inn}{{\mathrm{Inn}}} \newcommand{\Spf}{{\mathrm{Spf}}}
\newcommand{\Mat}{{\mathrm{Mat}}}


    \font\cyr=wncyr10  \newcommand{\Sha}{\hbox{\cyr X}}
    \newcommand{\wt}{\widetilde} \newcommand{\wh}{\widehat} \newcommand{\ck}{\check}
    \newcommand{\pp}{\frac{\partial\bar\partial}{\pi i}}
    \newcommand{\pair}[1]{\langle {#1} \rangle}
    \newcommand{\wpair}[1]{\left\{{#1}\right\}}
    \newcommand{\intn}[1]{\left( {#1} \right)}
    \newcommand{\norm}[1]{\|{#1}\|}
    \newcommand{\sfrac}[2]{\left( \frac {#1}{#2}\right)}
    \newcommand{\ds}{\displaystyle}
    \newcommand{\ov}{\overline}
    \newcommand{\Gros}{Gr\"{o}ssencharaktere}
    \newcommand{\incl}{\hookrightarrow}
    \newcommand{\lra}{\longrightarrow}
     \newcommand{\ra}{\rightarrow}
    \newcommand{\imp}{\Longrightarrow}
    \newcommand{\lto}{\longmapsto}
    \newcommand{\bs}{\backslash}
    \newcommand{\nequiv}{\equiv\hspace{-7.8pt}/}
    \theoremstyle{plain}


\definecolor{energy}{RGB}{114,0,172}
\definecolor{freq}{RGB}{45,177,93}
\definecolor{spin}{RGB}{251,0,29}
\definecolor{signal}{RGB}{203,23,206}
\definecolor{circle}{RGB}{217,86,16}
\definecolor{average}{RGB}{203,23,206}
\newcommand{\K}{\operatornamewithlimits{K}}
\colorlet{shadecolor}{gray!20}
\pgfplotsset{compat=1.9}
\def\N{10}
\def\M{4}
\usepgflibrary{fpu}

\algnewcommand\algorithmicinput{\textbf{Input:}}
\algnewcommand\algorithmicoutput{\textbf{Output:}}

\algnewcommand\Input{\item[\algorithmicinput]}
\algnewcommand\Output{\item[\algorithmicoutput]}

\def\upint{\mathchoice%
    {\mkern13mu\overline{\vphantom{\intop}\mkern7mu}\mkern-20mu}%
    {\mkern7mu\overline{\vphantom{\intop}\mkern7mu}\mkern-14mu}%
    {\mkern7mu\overline{\vphantom{\intop}\mkern7mu}\mkern-14mu}%
    {\mkern7mu\overline{\vphantom{\intop}\mkern7mu}\mkern-14mu}%
  \int}
\def\lowint{\mkern3mu\underline{\vphantom{\intop}\mkern7mu}\mkern-10mu\int}



\makeatletter
\let\c@equation\c@thm
\raggedbottom
\makeatother
\numberwithin{equation}{section}
%--------Meta Data: Fill in your info------
\author{Hojune Lee, 20210541}

\title{Cryptography}
\begin{document}

\begin{titlepage}
    \begin{center}
        \vspace*{9cm}
            
        \Huge
        \textbf{Cryptography}
    
        \vspace{1cm}
        \large
        Summary note, Lectured by Zhandry
        \vspace{3cm}
        
        \LARGE
        \textbf{Hojune Lee}
            
        \vspace{8cm}
            
        \normalsize
        \textbf{School of Computing, KAIST}\\  
    \end{center}
\end{titlepage}

\hypersetup{linkcolor=black}
\tableofcontents
\hypersetup{linkcolor=blue}

\newpage 

\chapter{Introduction to Cryptography}

\section{Encryption}

There will be many notational things and remarks and pre-requested things. 
Since this is just \lq summary note of summary note(or more depth)', many of things will be assumed. 
For here, I'm familiar with basic undergraduated mathematics and computation theory(e.g., Turing machine) 
so I'll omit such details. Let's begin. 

\begin{tcolorbox}[colback=yellow!10!white,colframe=red!75!black,title=Encryption(Symmetric key or Secret Key)]

(Symmetric key or Secret key) Encryption scheme is pair (\textsc{Enc}, \textsc{Dec}) of two (probabilistic) polynomial time algorithms such that 
\[\textsc{Enc} : (\text{key, plain text}) \mapsto \text{cipher text}\]
\[\textsc{Dec} : (\text{key, cipher text}) \mapsto \text{plain text}\]

In symmetric key encryption scheme, assume that both side users have common key $k$. 
Then what we want to guarantee is written by: 
\[\Pr\left[\textsc{Dec}(k, \textsc{Enc}(k, m)) = m, \underbrace{k \xleftarrow{\mathdollar} \{0,1\}^\lambda}_{\text{environment}} \right] = 1\] 

Means that, under the setting that $k$ is purely random among $\{0,1\}^\lambda$ where $\lambda$ is key length, 
above event happens almost always. 

\end{tcolorbox}

In Cryptography, we interested about security of such schemes. Since the word \lq security' is quite abstracted, 
there are bunch of things to define. First, we must define our problem (or settings). 
In security problem, we always assume there exists unknown adversary $A$. Then the settings here can be: 
\begin{itemize}
    \item Could $A$ obtain the full cipher text?
    \item What's the goal of $A$? Obtain full plain text, or some function if the plain text?
    \item Could $A$ know the \textsc{Enc, Dec} and exploit them? 
    \item ...
\end{itemize}

Our goal is analyze under these various assumptions. Let's see the first example. 

\begin{tcolorbox}[colback=yellow!10!white,colframe=blue!75!white,title=Indistinguishability Chosen Plaintext Attck]
$\textsc{IND-CPA-EXP}_b(A, \lambda)$

Let $A$ be an adversary, $C$ be a challenger. $b$ is parameter bit, here we assume that $A$ never access to $b$. 
\begin{enumerate}
    \item $C$ chooses purely(uniformly) random key $k \xleftarrow{\mathdollar} \{0,1\}^\lambda$
    \item $A$ can choose any 2 plain text message $m_0, m_1$ then send both of them to $C$. 
    \item $C$ select $m_b$ and then construct $c \gets \textsc{Enc}(k, m_b)$ then send it to $A$. 
    \item $A$ can repeat above step arbitrary. Each time of repitition, we count $A$ one unit of time. 
    \item Finally, $A$ guess $b'$ for $b$. 
\end{enumerate}

\end{tcolorbox}

Intuitively, $b$ is sort of information that $C$ hiding. If $A$ can guess correctly with high probability, the encryption scheme reveal 
few information about $C$'s secret to $A$. So, the ideal security under this setting can be defined by following.

\begin{tcolorbox}[colback=yellow!10!white,colframe=brown!75!black,title=Indistinguishable under Chosen Plaintext Attack]

An encryption scheme $(\textsc{Enc}, \textsc{Dec})$ is \textsc{IND-CPA} secure(indistinguishable under a chosen plaintext attack) 
when for all adversary $A$ who have polynomial resources limitation, 
\[\bigg| \Pr\left[1 \gets \textsc{IND-CPA-EXP}_0(A, \lambda)\right] - \Pr\left[1 \gets \textsc{IND-CPA-EXP}_1(A, \lambda)\right]\bigg| < \varepsilon(\lambda)\]

Means that, two probability that $A$ is wrong or correct is almost equal. 

\end{tcolorbox}

We often call such scheme by \textsc{CPA} secure. However, deterministic \textsc{Enc} can not be \textsc{CPA} secure. 
$A$ can send $(m_0, m_0)$ and $(m_1, m_1)$. Then $A$ obtain $c_0, c_1$ which are deterministic w.r.t. $m_0, m_1$. Then $A$ can send $(m_0, m_1)$ then 
can obtain correct $b$ with probability 1. 

\newpage 

\section{Pseudo Random Functions}

Since we've seen that deterministic is not suitable for security. So we need any sort of notion of Random. 
Here, we define what is computable random(pseudo random), and pure random. 

\begin{tcolorbox}[colback=yellow!10!white,colframe=red!75!black,title=Pseudo Random Function]

Pseudo Random Function(\textsc{PRF}) is a deterministic polynomial time computable function such that 
\[\textsc{PRF}: \{0,1\}^\lambda \times \{0,1\}^{n(\lambda)} \rightarrow \{0,1\}^{m(\lambda)}\]

\end{tcolorbox}

It can be understood by \textsc{PRF} maps $(\text{key}, \text{seed}) \mapsto \text{Random string}$. 
To it work as valid and secure random function, we do similar thing to previous section. Let's define following game. 

\begin{tcolorbox}[colback=yellow!10!white,colframe=blue!75!white,title=Pseudo Random Function Experiments]

$\textsc{PRF-EXP}_b(A, \lambda)$

Let $A$ be an \textsc{PPT} adversary, $C$ be a challenger. $b$ is parameter bit, here we assume that $A$ never access to $b$. 

\begin{enumerate}
    \item If $b = 0$, $C$ choose purely random key $k \xleftarrow{\mathdollar} \{0,1\}^\lambda$. If $b = 1$, $C$ initializes an empty list $L$.
    \item $A$ choose any seed $x \in \{0,1\}^{n(\lambda)}$ then send it to $C$. Then, $C$ responds as follow: 
    \begin{itemize}
        \item If $b = 0$, $C$ return $y \gets \textsc{PRF}(k, x)$
        \item If $b = 1$, if $(x, y) \in L$ then return $y$ else generate pure random $y$ and $L \gets L :: (x, y)$
    \end{itemize}
    \item $A$ can repeat this process arbitrary and one repition charge one time unit.
    \item Finally, $A$ guess $b'$ for $b$. 
\end{enumerate}

\end{tcolorbox}

In other words, the goal of $A$ is distinguish computed random and pure random in polynomial time. If 
it is possible, then such \textsc{PRF} is not so good. So, we define secure \textsc{PRF} via above experiment. 

\begin{tcolorbox}[colback=yellow!10!white,colframe=brown!75!black,title=PRF secure]

A \textsc{PRF} is called PRF secure if for any PPT adversary $A$, 
\[\bigg| \Pr\left[1 \gets \textsc{PRF-EXP}_0(A, \lambda)\right] - \Pr\left[1 \gets \textsc{PRF-EXP}_1(A, \lambda)\right] \bigg| < \varepsilon(\lambda)\]

\end{tcolorbox}

\newpage 

\section{CPA-secure secret key encryption wirh PRF}

Then we can construct encryption scheme with such PRF. Let's see the following construction. 

\begin{tcolorbox}[colback=yellow!10!white,colframe=blue!75!white,title=Encryption with PRF]

Let's assume that \textsc{PRF} is given.
\begin{itemize} 
\item \textsc{Enc}$(k, m)$: Given key $k \in \{0,1\}^\lambda$, plain text $m \in \{0,1\}^{m(\lambda)}$. 
Choose $r \in \{0,1\}^{n(\lambda)}$ randomly (pure or pseudo) then compute $c \gets \textsc{PRF}(k,r) \oplus m$
Output of encryption is $(r, c)$. 
\item \textsc{Dec}$(k, (r,c))$: Given key $k \in \{0,1\}^{\lambda}$, cipher text $c \in \{0,1\}^{m(\lambda)}$ and chosen seed $r \in \{0,1\}^{n(\lambda)}$.
Compute $m' \gets \textsc{PRF}(k, r) \oplus c = \textsc{PRF}(k,r) \oplus (\textsc{PRF}(k, r) \oplus m) = m$ 
\end{itemize}
\end{tcolorbox}

Then what we want to show is that above encryption is CPA-secure. 

\begin{tcolorbox}[colback=yellow!10!white,colframe=brown!75!black,title=Encryption with PRF is CPA-secure]

If given \textsc{PRF} is \textsc{PRF}-secure and $2^{n(\lambda)}$ is super-polynomial, 
then above $(\textsc{Enc}, \textsc{Dec})$ is CPA-secure. 

\end{tcolorbox}

\begin{proof}
    
Suppose not. Then since $(\textsc{Enc}, \textsc{Dec})$ is not CPA-secure, we can write 
\[\bigg| \Pr\left[1 \gets \textsc{IND-CPA-EXP}_0(A, \lambda)\right] - \Pr\left[1 \gets \textsc{IND-CPA-EXP}_1(A, \lambda)\right] \bigg| \geq \varepsilon(\lambda)\]
Our goal is breaking the \textsc{PRF}-security of given pseudo random function with this inequality. 
For intuitive idea, our goal is construct following. 
\begin{enumerate}
    \item We'll define 2 more games, called $H_1, H_2$. Let $\textsc{IND-CPA-EXP}_0 = H_1$ and 
    
    $\textsc{IND-CPA-EXP}_1 = H_3$. 
    \item Then, we already have 
    \[\bigg| \Pr[1 \gets H_0] - \Pr[1 \gets H_3] \bigg| \geq \varepsilon(\lambda)\]
    \item Then, we can obtain 
    \[\bigg| \underbrace{\Pr[1 \gets H_0] - \Pr[1 \gets H_1]}_A + \underbrace{...}_B + \underbrace{\Pr[1 \gets H_2] - \Pr[1 \gets H_3]}_C \bigg| \geq \varepsilon(\lambda)\]
    \item Then, $|A| + |B| + |C| \geq |A+B+C| \geq \varepsilon(\lambda)$, implies that 
    \[\exists i \in \{0,1,2\}, \bigg| \Pr\left[1 \gets H_i\right] - \Pr\left[1 \gets H_{i+1}\right] \bigg| \geq \frac{1}{3}\varepsilon(\lambda) \]
    \item Finally, all each cases contradicts to \textsc{PRF}-security. 
\end{enumerate}

Then what we need to do is that, define $H_1, H_2$ and show all cases $i \in \{0, 1, 2\}$ fails. 

\begin{itemize}
    \item $H_0$ : $A$ sends $m_0, m_1$ to $C$. Then $C$ choose $m_0$, and choose random $r \in \{0, 1\}^{n(\lambda)}$, 
    then send $A$ back to $(r, m_0 \oplus \textsc{PRF}(k, r))$
    \item $H_1$ : $C$ initialize $L = \emptyset$. $A$ sends $m_0, m_1$ to $C$. 
    $C$ choose random $r \in \{0,1\}^{n(\lambda)}$ then search $(r, y)$ in $L$. If not exists, then 
    generate pure random $y$ and $L \gets L::(r, y)$ and return $(r, m_0 \oplus y)$
    \item $H_2$ : $C$ initialize $L = \emptyset$. $A$ sends $m_0, m_1$ to $C$. 
    $C$ choose random $r \in \{0,1\}^{n(\lambda)}$ then search $(r, y)$ in $L$. If not exists, then 
    generate pure random $y$ and $L \gets L::(r, y)$ and return $(r, m_1 \oplus y)$
    \item $H_3$ : $A$ sends $m_0, m_1$ to $C$. Then $C$ choose $m_1$, and choose random $r \in \{0, 1\}^{n(\lambda)}$, 
    then send $A$ back to $(r, m_1 \oplus \textsc{PRF}(k, r))$
\end{itemize}

All of these are quite variance of \textsc{IND-CPA-EXP} game. Now, our remaining goal is that 
if one of $H_i, H_{i+1}$ is distinguishable, then it contradicts to \textsc{PRF}-security. 

\begin{enumerate}
    \item $(i = 0,2)$ Here, suppose that $A$ has power:
    \[\bigg| \Pr[1 \gets H_0] - \Pr[1 \gets H_1] \bigg| \geq \frac{1}{3}\varepsilon(\lambda)\]$A$ can break the \textsc{PRF}-security of given PRF. 
    Now $B$ join the $\textsc{PRF-EXP}_b(A, \lambda)$. Then, $A$ prepare $(m_0, m_1)$ and send it to $B$. 
    $B$ generate random string $r \in \{0,1\}^{n(\lambda)}$ and send $B$'s challenger. If $b = 0$, it gives pseudo random $y$, or $b = 1$, it gives pure random $y$. 
    Anyway, $B$ make $(r, y \oplus m_0)$ then give it back to $A$. Then $A$ will distinguish $y$ is pseudo random or pure random with probability $\varepsilon(\lambda)/3$. 
    Then, $B$ takes the output of $A$ and submit it to its challenger. Then $B$ also distinguish pseudo random and pure random with probability $\varepsilon(\lambda)/3$, contradiction. 
    Vice versa for $i = 2$.
    \item $(i = 1)$ Here, suppose that $A$ has power: 
    \[\bigg| \Pr[1 \gets H_1] - \Pr[1 \gets H_2] \bigg| \geq \frac{1}{3}\varepsilon(\lambda)\]
    Here, note that $A$ can distinguish $m_0 \oplus y$ and $m_1 \oplus y$ with probability $\varepsilon(\lambda)/3$. 
    However, $y$ is pure random string with $m(\lambda)$ length. Since $y$ is purely random, 
    $m_0 \oplus y, m_1 \oplus y$ both are also purely random strings. So it is generally impossible to distinguish them. 
    One possible case is that, $r$ is small so there can be collison in $y$. By the mechanism, 
    we can possibily obtain $2^{n(\lambda)}$ amount of $y$ so the collision probability is $2^{-n(\lambda)}$. Since $2^{n(\lambda)}$ is super polynomial, 
    this probability is negligible. $A$ can generate $q$ queries, collision probability is $q^2/2^{n(\lambda)}$ which is negligible. So $A$ can't distinguish $H_1, H_2$, this case is impossible. 
\end{enumerate}

\end{proof}

\chapter{Cryptography From Minimal Assumptions}

In the previous chapter, we've seen that how can construct basic encryption via secure PRF. 
By the way, we just assume that PRF is just a mapping from $\{0,1\}^\lambda \times \{0,1\}^{n(\lambda)}$ 
into $\{0,1\}^{m(\lambda)}$ which satisfy PRF security: 
\[\left|\Pr[1 \gets \text{PRF-EXP}_0(A, \lambda)] - \Pr[1 \gets \text{PRF-EXP}_1(A, \lambda)]\right| < \varepsilon(\lambda)\]
Here $1 \gets \text{PRF-EXP}_b(A, \lambda)$ means that $A$ decide that itself gets some value via purely random manner. 
Natural wondering is that how PRF works. 

\section{Pseudo Random Generators}

\begin{tcolorbox}[colback=yellow!10!white,colframe=brown!75!black,title=Pseudo Random Generator]

PRG $G$ is deterministic polynomial time computable function such that 
\[G: \{0, 1\}^{\lambda} \rightarrow \{0, 1\}^{\lambda + s(\lambda)}, \quad s(\lambda) \geq 1\]

\end{tcolorbox}

\begin{tcolorbox}[colback=yellow!10!white,colframe=blue!75!white,title=Secure PRG]

Let $A$ be a PPT adversary who tries to determine whether given string in $\{0,1\}^{\lambda + s(\lambda)}$ is 
randomly generated or deterministically generated. 
\[A(q \in \{0,1\}^{\lambda + s(\lambda)}) = \begin{cases}
    0 & \text{A think } q \xleftarrow{\$} \{0, 1\}^{\lambda + s(\lambda)} \\
    1 & \text{A think }\exists G: \{0,1\}^\lambda \rightarrow \{0, 1\}^{\lambda + s(\lambda)}, q = G(x)
\end{cases}\]

Then, we say that PRF $G$ is secure when following inequality holds for negligible function $\varepsilon$: 
\[\left|\Pr[A(G(x)) = 1 ~:~ x \xleftarrow{\$} \{0, 1\}^\lambda] - \Pr[A(y) = 1 ~:~ y \xleftarrow{\$} \{0,1\}^{\lambda + s(\lambda)}]\right| < \varepsilon(\lambda)\]

\end{tcolorbox}

This means that, for secure PRF $G$, any PPT adversary can't determine whether his query is 
randomly generated or deterministically generated by some $G$. 

\begin{tcolorbox}[colback=yellow!10!white,colframe=gray!75!black,title=Remark?]

If $A$ has exponential power, than try $\CO(2^\lambda)$ times and if $A$ saw quite amount of collisions, 
then determine that query is deterministically generated. 
This probability is non-negligible. 

\end{tcolorbox}

As you can see from the codomain of PRG, one deterministic way to generate $s(\lambda)$ bits longer 
string from $\{0,1\}^\lambda$ might be precise inductive extension. 
Let's define 
\[G: \{0, 1\}^\lambda \rightarrow \{0,1\}^{\lambda + 1}\]
works with deterministic mappings $G_0: \{0,1\}^{\lambda} \rightarrow \{0,1\}^{\lambda}, b: \{0,1\}^{\lambda} \rightarrow \{0,1\}$.
\[G: x \mapsto (G_0(x), b(x))\]
Then we can define inductively extension $G_s: \{0,1\}^\lambda \rightarrow \{0,1\}^{\lambda + s}$ via: 
\[G_s: x \mapsto (x_s, o_1, \cdots, o_s)\]
where 
\[x_0 = x, \quad x_{i} = G_0(x_{i-1}), \quad o_i = b(x_{i-1})\]
In other words, $G_s$ is inductive chain of one-bit extension. We can prove that if such one-bit extension is 
secure PRG, then $G_s$ also. 

\begin{tcolorbox}[colback=yellow!10!white,colframe=red!75!black,title=PRG Security Theorem]

If one bit extension $G$ is secure PRG, then $G_s$ is also secure PRG for any $s \in \BN$. 

\end{tcolorbox}

\begin{proof}
    
We'll use Hybrid Experiment Scheme here also. Let's denote the working of $G_s$ in this: 
\[\{0, 1\}^\lambda \xrightarrow{\$} x_0 \rightarrow (x_1 = G_0(x_0), o_1 = b(x_0)) \rightarrow \cdots \rightarrow (x_s, o_1, \cdots, o_s) \]
Then, we can imagine slight variance of this construction for $1 \leq i \leq s$.
\[\underbrace{\left(\{0, 1\}^\lambda \xrightarrow{\$} x_i, \{0, 1\} \xrightarrow{\$} o_1, \cdots, o_i\right)}_{\text{Intercept Step}} \rightarrow (x_{i+1} = G_0(x_i), o_{i+1} = b(x_i)) \rightarrow \cdots \rightarrow (x_s, o_1, \cdots, o_s)\]
We call each of them are $\textsc{Hybrid}_i$. Here, $1 \gets \textsc{Hybrid}_i(A, y)$ is $A$ decide this string $y$ is generated via $\textsc{Hybrid}_i$ method. Note that $\textsc{Hybrid}_s$ is uniformly random choice of entire $(x_s, o_1, \cdots, o_s) \in \{0, 1\}^{\lambda + s}$. 
And denote original $G_s$ be $\textsc{Hybrid}_0$. 

Now, suppose that $G_s$ is non-secure PRG. Then $\textsc{Hybrid}_0$ and $\textsc{Hybrid}_s$ is distinguishable and following holds: 
\[\left|\Pr[A(G_s(x)) = 1 ~:~ x \xleftarrow{\$} \{0, 1\}^\lambda] - \Pr[A(y) = 1 ~:~ y \xleftarrow{\$} \{0,1\}^{\lambda + s(\lambda)}]\right| \geq \varepsilon(\lambda)\]
However, we can denote this into what we've defined: 
\begin{align*}
    & \left|\Pr[A(G_s(x)) = 1 ~:~ x \xleftarrow{\$} \{0, 1\}^\lambda] - \Pr[A(y) = 1 ~:~ y \xleftarrow{\$} \{0,1\}^{\lambda + s(\lambda)}]\right| \geq \varepsilon(\lambda) \\
    & \implies \big|\Pr\left[1 \gets \textsc{Hybrid}_0(A, y)\right] - \Pr\left[1 \gets \textsc{Hybrid}_s(A, y)\right]\big| \geq \varepsilon(\lambda) \\
    & \implies \left|\sum_{i = 0}^{s-1} \Pr\left[1 \gets \textsc{Hybrid}_{i+1}(A, y)\right] - \Pr\left[1 \gets \textsc{Hybrid}_i(A, y)\right]\right| \geq \varepsilon(\lambda) \\
    & \implies \text{for some $i$, } \left|\Pr\left[1 \gets \textsc{Hybrid}_{i+1}(A, y)\right] - \Pr\left[1 \gets \textsc{Hybrid}_i(A, y)\right]\right| \geq \frac{\varepsilon(\lambda)}{s}
\end{align*}

With this, now we'll show that this breaks the PRG secureness of $G$. Let $B$ be another 
adversary for breaking $G$. Suppose that $B$ gets $y \in \{0, 1\}^{\lambda + 1}$. Goal of $B$ is determining whether 
$y \xleftarrow{\$} \{0, 1\}^{\lambda + 1}$ or $y = G(x), x \xleftarrow{\$} \{0, 1\}^\lambda$. 
In perspective of $B$, let 
\[y = (x_{i+1}, o_{i+1}), \quad x_{i+1} \in \{0, 1\}^{\lambda}, o_{i+1} \in \{0, 1\}\]
Now $B$ choose $i$ bits for: 
\[\text{$B$ choose : } o_1, \cdots, o_i \xleftarrow{\$} \{0, 1\}\]
And since $B$ knows and access to $G$, $B$ can compute 
\[x_{i+2} = G_0(x_{i+1}), o_{i+2} = b(x_{i+1}) \rightarrow \cdots \rightarrow x_{s} = G_0(x_{s-1}), o_s = b(x_{s-1})\]
So now $B$ obtain 
\[z = (\underbrace{x_s}_{\text{computed}}, \underbrace{o_1, \cdots, o_i}_{\text{uniform random}}, \underbrace{o_{i+1}}_{\text{query}}, \underbrace{o_{i+2}, \cdots, o_{s}}_{\text{computed}})\]
$B$ send this for $A$. Since $A$ can distinguish $\textsc{Hybrid}_i$ and $\textsc{Hybrid}_{i+1}$, 
when $A$ see $z$, 
\begin{align*}
    \frac{\varepsilon(\lambda)}{s} & \leq \left|\Pr\left[1 \gets \textsc{Hybrid}_{i+1}(A, z)\right] - \Pr\left[1 \gets \textsc{Hybrid}_i(A, z)\right]\right| \\
    & =  \bigg|\Pr[A \text{ decide $z$ is generated by $\textsc{Hybrid}_{i+1}$}] \\
    & \qquad \qquad - \Pr[A \text{ decide $z$ is generated by $\textsc{Hybrid}_{i}$}] \bigg| \\
    &= \bigg|\Pr[A \text{ decide $x_{i+1}, o_{i+1}$ is randomly chosen}] \\
    & \qquad \qquad - \Pr[A \text{ decide $x_{i+1}, o_{i+1}$ is computed via $G$}] \bigg| \\
\end{align*}
Now $B$ catch output of $A$, then $B$ also decide $y = (x_{i+1}, o_{i+1})$ is randomly generated or 
computed bia $G$ with probability larger than $\varepsilon(\lambda)/s$. So, 
\[\left|\Pr[B(G(x)) = 1 ~:~ x \xleftarrow{\$} \{0, 1\}^\lambda] - \Pr[B(y) = 1 ~:~ y \xleftarrow{\$} \{0, 1\}^{\lambda + 1}]\right| \geq \frac{\varepsilon(\lambda)}{s}\]
So this is contradiction to $G$ is secure PRF. So, $G_s$ is also secure PRF. 
\end{proof}

Then, now we'll see the example of construction of PRF from PRG. Choose one key: 
\[k \in \{0, 1\}^{\lambda}\]
Then suppose that our Secure Pseudo Random Generator $G$ be length double-expansion function: 
\[G: \{0, 1\}^\lambda \rightarrow \{0, 1\}^{2\lambda}\]
With this PRG, we can generate $2^{n(\lambda)} \cdot \lambda$ bits string: 
\[
\begin{tikzcd}
 & & & k \arrow[d, "G"] & & & \\
 & & & G(k) \arrow[lld] \arrow[rrd] & & & \\
 & L^1_0 \arrow[d] & & & & L^1_1 \arrow[d] & \\
 & G(L^1_0) \arrow[ld] \arrow[rd] & & & & G(L^1_1) \arrow[ld] \arrow[rd] & \\
 L^2_{00} & & L^2_{01} & & L^2_{10} & & L^2_{11} \\
 & & & \vdots & & & 
\end{tikzcd}
\]
So, in depth $n(\lambda)$, we can make $2^{n(\lambda)}$ of $\lambda$-length bits strings. 
With this, our PRF works: 
\[\textsc{PRF}: \{0, 1\}^\lambda \times \{0, 1\}^{n(\lambda)} \rightarrow \{0, 1\}^\lambda\]
Explicitly, 
\[\textsc{PRF}: (k, s) \mapsto L^{n(\lambda)}_s \in \{0, 1\}^\lambda\]
Although the entire space of co-domain is $2^\lambda$, when $k$ is fixed, the number of leaves is 
$2^{n(\lambda)}$. However we can compute efficiently when we know $s$. For each $i$-th branch, 
if $s_i = 0$ then choose left and $s_i = 1$ then choose right of given $2\lambda$ length output of $G$. 

\begin{tcolorbox}[colback=yellow!10!white,colframe=red!75!black,title=Secure PRG makes Secure PRF]

When $G$ is secure PRG, then above PRF is secure PRF. 

\end{tcolorbox}

\begin{tcolorbox}[colback=yellow!10!white,colframe=gray!75!black,title=Remark]

We can't apply Hybrid Experiment method here directly sicne we need $2^{n(\lambda)}$ hybrid experiments, 
and then $\varepsilon' > \varepsilon / 2^{n(\lambda)}$ also can be negligible. 

\end{tcolorbox}

Note that we can compute specific leaf in polynomial time as mentioned. 
Thus the time complexity to compute PRF value is $\CO(n(\lambda) \cdot \text{time}(G))$ so the running time of 
PRF is polynomial. Now let's prove that if $G$ is secure PRG, than above PRF is also secure. 

\begin{tcolorbox}[colback=yellow!10!white,colframe=blue!75!black,title=Length Doubling PRF Construction is Secure]

\textit{If $G$ is secure PRG, then above PRF construction is secure.}

\end{tcolorbox}

\begin{proof}
    
Suppose PRF is not secure, then adversary $A$ can obtain following advantage: 
\[\bigg|\Pr\left[1 \gets \text{PRF-EXP}_0(A, \lambda)\right] - \Pr\left[1 \gets \text{PRF-EXP}_1(A, \lambda)\right]\bigg| \geq \varepsilon(\lambda)\]
Now we'll define hybrid game $H_0, \cdots, H_n$ by following way. 
\[H_i := \begin{cases} 
    \text{Game with Original PRF} & i = 0 \\
    \text{Replace $L^1, \cdots, L^i$ into uniform random string} & 1 \leq i \leq n
\end{cases}\]
In other words, we can view $H_i$ as $n - i$ times application of PRG $G$ mechanically.
Then we can rewrite $\text{PRF-EXP}_0(A, \lambda)$ as $H_0$ and $\text{PRF-EXP}_1(A, \lambda)$ as $H_n$. So our assumption implies that $A$ can distinct $H_0$ and $H_n$, so by triangular inequality, 
there exists $i \in [n]$ such that 
\[\bigg|\Pr\left[1 \gets H_{i-1}(A)\right] - \Pr\left[1 \gets H_{i}(A)\right]\bigg| \geq \frac{\varepsilon(\lambda)}{n}\]
So $A$ can distinct each $2^i$ of $\lambda$-length string is randomly generated or generated from $i-1$ step. 
$2^{i-1}$ PRG is participated here so we can't use hybrid method directly. The key idea is that: $A$ is polynomial adversary, so $A$ can distinguish $H_{i-1}$ and $H_i$ with at most polynomial queries. 
Here exists 2 possible cases and we'll show how to prove for each case. 

\begin{enumerate}
    \item When $A$ first commits total $q(\lambda)$ queries and then send it to challenger. 
    \item When $A$ send total $q(\lambda)$ queries without pre-commits. 
\end{enumerate}
In the first case, we can know at most $q(\lambda)$ of $L^{i-1}$ nodes that $A$ will touch during challenging. 
Let's denote such strings by 
\[L^{i-1}_{1}, \cdots, L^{i-1}_{m}, \quad m \leq q(\lambda)\]
Construct $m$ hybrid games $H_i^0, \cdots, H_i^{m}$ by $H_i^j$ : childs of $L^{i-1}_1, \cdots, L^{i-1}_j$ is generated by PRG $G$ and remainings are purely random. Then, $H_{i-1}$ is same with $H_i^m$ and $H_i$ is same with $H^{0}_i$. 
So, $A$ can distinguish $H_i^j, H_i^{j-1}$ for some $j \in [m]$ 
\[\bigg|\Pr\left[1 \gets H^{j}_i(A)\right] - \Pr\left[1 \gets H^{j-1}_i(A)\right]\bigg| \geq \frac{\varepsilon(\lambda)}{n \cdot m} \geq \frac{\varepsilon(\lambda)}{n \cdot q(\lambda)}\]
So, $A$ can distinct $H_i^j$ and $H_i^{j+1}$ efficiently. Then the PRG $G$ adversary $B$ can exploit $A$ to break PRG security of $G$. 
$B$ gets $Z \gets \{0,1\}^{2\lambda}$ and goal is distinguish $Z$ is from pure random generation or $G$. Now, $B$ construct PRF simulation by child of $L^{i-1}_j$ be $Z$. 
Then $A$ can distinguish $Z$ non-negligibly, so obtaining 
\[\bigg|\Pr[1 \gets B(y), y \xleftarrow{\$} \{0,1\}^{2\lambda}] - \Pr[1 \gets B(G(x)), x \xleftarrow{\$} \{0,1\}^\lambda]\bigg| \geq \frac{\varepsilon(\lambda)}{n \cdot q(\lambda)}\]
Second condition is called Adaptive Queries condition and need more tricks. 
We'll define $L^{i-1}_1, \cdots, L^{i-1}_m$ where $L^{i-1}_j$ means that $j$-th new arrival of level $i-1$ during simulation of $A$. We don't need to exact spot that $A$ distinct, but we can assume that $A$ can distinct when $k$-th new arrival of $i-1$ level. 
Under this, we can simulate $B$ as similar with above and obtaining $B$ breaks security of $G$, contradiction. 

\end{proof}

\section{One Way Functions and One Way Permutations}

We studied PRG to construct PRF. However, there are also another mistery: How can we implement PRG? 
In this section, we define one-way functions and permutations and then implement PRG. 

\begin{tcolorbox}[colback=yellow!10!white,colframe=brown!75!black,title=One Way Function]

A \textit{one-way function} is a deterministic poly-time function $f: \{0,1\}^\lambda \rightarrow \{0,1\}^{n(\lambda)}$ 
such that for any PPT adversary $A$ exists negligible $\varepsilon$ such that 
\[\Pr[f(A(f(x))) = f(x), x \xleftarrow{\$} \{0,1\}^\lambda] < \varepsilon(\lambda)\]
Means that, PPT adversary $A$ tries to find pre-image of $f(x)$ such that randomly chosen but success probability is negligible. 
We call that $f$ is \textit{one-way permutation} if $f$ is bijective and $n(\lambda) = \lambda$. In this case, 
$A$ tries to find unique pre-image of $f(x)$. 
\end{tcolorbox}

\begin{tcolorbox}[colback=yellow!10!white,colframe=orange!75!white,title=Goal of This Section]

Implement PRG via one-way permutation(OWP) and \lq hardcore bit'. 

\end{tcolorbox}

\begin{tcolorbox}[colback=yellow!10!white,colframe=brown!75!black,title=Hardcore bit]

Let $f: X \rightarrow Y$ be a one-way function. A hardcore bit for $f$ is a function $hc_f: X \rightarrow \{0,1\}$ 
such that for all PPT adversary $A$ there exists $\varepsilon$ negligible such that 
\[\bigg|\Pr[1 \gets A(f(x), hc_f(x)) : x \gets X] - \Pr[1 \gets A(f(x), b) : x \gets X, b \gets \{0,1\}]\bigg| < \varepsilon(\lambda)\]
That means that, $f$ is (intuitively) independent with such one bit. $f$ does not be any hint to determine this bit is 
generated from purely random or hardcore bit function. 
\end{tcolorbox}

We already see the basic building blocks of PRG, which was $G_0$ and $b$. We'll explain this explicitly via OWP and hardcore bit. 
Let $f$ be one-way permutation and $hc$ is its hardcore bit function. 
Then $G(x) = (f(x), hc(x))$ becomes one-bit stretch PRG. This PRG is secure. Since $f$ is one-way permutation, if $x$ is pure randomly chosen, 
we can assume that $f(x)$ is also purely random bits string. So that, $A$ can't distinct $f(x)$ and purely random $\lambda$ bits string, and 
last hardcore bit also. 

\begin{tcolorbox}[colback=yellow!10!white,colframe=gray!75!black,title=Remarks]

Are there any $hc$ function such that hardcore for any OWP? Answer is no. Let $f$ be any OWF. Then we construct 
new OWF $f'(x) = f(x) ~||~ hc(x)$. Then adversary easily know that the last bit of $f'(x)$ is hardcore bit. So, $hc$ is not hardcore bit for $f'$. 

\end{tcolorbox}
Now we state the main important theorem, to claim that every OWF has proper hardcore bit function. (So that for any OWF $f$, we can find $hc_f$ and construct PRG by $G(x) = (f(x), hc_f(x))$.) 

\begin{tcolorbox}[colback=yellow!10!white,colframe=blue!75!black,title=Goldreigh-Levin Theorem]

Let $F: \{0,1\}^n \rightarrow Y$ be a OWF. Define $F': \{0,1\}^{2n} \rightarrow \{0,1\}^n \times Y$ such that 
$F'(x, r) = (r, F(x))$. Then, $\langle x, r \rangle = \sum x_i r_i \pmod{2}$ is a hardcore for $F'$. 

\end{tcolorbox}

\begin{proof}
    
Suppose not. Then, some PPT adversary $A$ can break hardcoreness of $\langle x, r \rangle$ and obtain following inequality: 
\[\bigg|\Pr[A(r, F(x), \langle x, r \rangle) = 1] - \Pr[A(r, F(x), b) = 1]\bigg| > \varepsilon(n)\]
Where $r \xleftarrow{\$} \{0,1\}^n$, $x \xleftarrow{\$} \{0,1\}^n$ and $b \xleftarrow{\$} \{0,1\}$. 
Without loss of generality, 
\[\Pr[A(r, F(x), \langle x, r \rangle) = 1] - \Pr[A(r, F(x), b) = 1] > \varepsilon(n)\]
Now let's construct new agent $B$. $B$ takes $r, F(x) \in \{0,1\}^n$ and pure-randomly choose $t \in \{0,1\}$. 
Then run $A(r, F(x), t)$. 
\[B(r, F(x)) = \begin{cases}
    t & A(r, F(x), t) = 1 \\
    1 - t & A(r, F(x), t) = 0
\end{cases}\]
Then $B$ guess the hardcore function of given instance $r, F(x)$. Since $A$ has $\varepsilon(n)$ advantage to determine whether given instance is from hardcore function or pure random, 
\[\Pr[B(r, F(x)) = \langle x, r \rangle] > \frac{1}{2} + \varepsilon(n)\]
This is because that, if $t = \langle x, r \rangle$ then $A(r, F(x), t) = 1$ has higher probability. 
If $1 - t = \langle x, r \rangle$, then $A(r, F(x), t) = 0$ has higher probability. So, $B$ can guess $\langle x, r \rangle$ by exploiting $A$ with $\varepsilon(n)$ advantage. 
With agent $B$, our goal is breaking the OWF property of $F$. 
If $\varepsilon(n) = \frac{1}{2}$, it implies that $B$ always answer correctly. Then by standard basis bit string $e_i$, we take $n$ queries of $B(e_i, x) = \langle x, e_i \rangle = x_i$ and 
obtain $x$. If $\varepsilon(n) = \frac{1}{2} - \frac{1}{2n}$, then choose $n$ random bit string $r_1, \cdots, r_n$ and compute 
$B(r_i, F(x))$. They correctly answer $\langle x, r_i \rangle$ with probability $1 - \frac{1}{2n}$, so 
total-correct probability is $\frac{1}{2}$. If they are all correct, we can observe $n$ linear system of $\langle x, r_i \rangle = b_i$. Solving this gives us total $x$ with probability $1/2$. 
Third, suppose that $\varepsilon(n) = \frac{1}{4} + \gamma$ where $\gamma$ is non-negligible. 
Then, we can denote 
\[\Pr[B(r, F(x)) = \langle x, r \rangle :~ x, r \xleftarrow{\$} \{0,1\}^n] = \frac{3}{4} + \gamma\]
Then we can view this as joint distribution of $x, r$, so let's define 
\[p(x) = \Pr[B(r, F(x)) = \langle x, r \rangle : r \xleftarrow{\$} \{0,1\}^n ]\]
By definition, 
\[\BE_{x \xleftarrow{\$} \{0,1\}^n}[p(x)] = \frac{3}{4} + \gamma\]
Let's denote $Y(x) = 1 - p(x)$ then we can rewrite above into 
\[\BE_{x \xleftarrow{\$} \{0,1\}^n}[Y(x)] = \frac{1}{4} - \gamma\]
$Y(x)$ denotes the probability that $B$ guess wrong answer for fixed $x$ and uniform $r$. So, 
if there exists $x \in \{0,1\}^n$ such that $Y(x) < \frac{1}{4} - \frac{\gamma}{2}$, we call such $x$ is \lq good'. If $Y(x) \geq \frac{1}{4} - \frac{\gamma}{2}$ 
then such $x$ is called \lq bad'. 

\begin{lem}(Markov's Inequality)
    For random variable $X: \Omega \rightarrow \BR$ and function $f: \BR \rightarrow \BR^+$, 
    \[P(f(X) \geq a) \leq \frac{\BE[f(X)]}{a}\]
\end{lem}

Then, 
\[\Pr[x \text{ is bad}] = \Pr[Y(x) \geq \frac{1}{4} - \frac{\gamma}{2}] \leq \frac{\BE[Y(x)]}{\frac{1}{4}-\frac{\gamma}{2}} = \frac{1 - 4\gamma}{1 - 2\gamma}\]
So, 
\[\Pr[x \text{ is good}] \geq 1 - \frac{1-4\gamma}{1-2\gamma} = \frac{2\gamma}{1 - 2\gamma}\geq 2\gamma\]
Means that we can choose good $x$ via sampling with non-negligible probability. Now assume that $x$ is good. Let 
\[H(r) = B(r, F(x))\]
And for every $i \in [n]$, choose $r \in \{0,1\}^n$ randomly then compute 
\[H(r) \oplus H(r \oplus e_i)\]
If $B$ guess correctly boss of them, the value is 
\[\langle x, r \rangle \oplus \langle x, r \oplus e_i \rangle = \langle x, e_i \rangle = x_i\]
So, 
\[\Pr[H(r) \oplus H(r \oplus e_i) = x_i] \geq \left(\frac{3}{4} + \gamma\right)^2 \geq \frac{1}{2} + \gamma\]
So, we have $\gamma$-advantage to compute $x_i$, so for good $x$, we have non-negligible advantage to compute $x$. 
\end{proof}

\end{document}